% Options for packages loaded elsewhere
\PassOptionsToPackage{unicode}{hyperref}
\PassOptionsToPackage{hyphens}{url}
\PassOptionsToPackage{dvipsnames,svgnames,x11names}{xcolor}
%
\documentclass[
  letterpaper,
  DIV=11,
  numbers=noendperiod]{scrartcl}

\usepackage{amsmath,amssymb}
\usepackage{iftex}
\ifPDFTeX
  \usepackage[T1]{fontenc}
  \usepackage[utf8]{inputenc}
  \usepackage{textcomp} % provide euro and other symbols
\else % if luatex or xetex
  \usepackage{unicode-math}
  \defaultfontfeatures{Scale=MatchLowercase}
  \defaultfontfeatures[\rmfamily]{Ligatures=TeX,Scale=1}
\fi
\usepackage{lmodern}
\ifPDFTeX\else  
    % xetex/luatex font selection
\fi
% Use upquote if available, for straight quotes in verbatim environments
\IfFileExists{upquote.sty}{\usepackage{upquote}}{}
\IfFileExists{microtype.sty}{% use microtype if available
  \usepackage[]{microtype}
  \UseMicrotypeSet[protrusion]{basicmath} % disable protrusion for tt fonts
}{}
\makeatletter
\@ifundefined{KOMAClassName}{% if non-KOMA class
  \IfFileExists{parskip.sty}{%
    \usepackage{parskip}
  }{% else
    \setlength{\parindent}{0pt}
    \setlength{\parskip}{6pt plus 2pt minus 1pt}}
}{% if KOMA class
  \KOMAoptions{parskip=half}}
\makeatother
\usepackage{xcolor}
\setlength{\emergencystretch}{3em} % prevent overfull lines
\setcounter{secnumdepth}{-\maxdimen} % remove section numbering
% Make \paragraph and \subparagraph free-standing
\makeatletter
\ifx\paragraph\undefined\else
  \let\oldparagraph\paragraph
  \renewcommand{\paragraph}{
    \@ifstar
      \xxxParagraphStar
      \xxxParagraphNoStar
  }
  \newcommand{\xxxParagraphStar}[1]{\oldparagraph*{#1}\mbox{}}
  \newcommand{\xxxParagraphNoStar}[1]{\oldparagraph{#1}\mbox{}}
\fi
\ifx\subparagraph\undefined\else
  \let\oldsubparagraph\subparagraph
  \renewcommand{\subparagraph}{
    \@ifstar
      \xxxSubParagraphStar
      \xxxSubParagraphNoStar
  }
  \newcommand{\xxxSubParagraphStar}[1]{\oldsubparagraph*{#1}\mbox{}}
  \newcommand{\xxxSubParagraphNoStar}[1]{\oldsubparagraph{#1}\mbox{}}
\fi
\makeatother

\usepackage{color}
\usepackage{fancyvrb}
\newcommand{\VerbBar}{|}
\newcommand{\VERB}{\Verb[commandchars=\\\{\}]}
\DefineVerbatimEnvironment{Highlighting}{Verbatim}{commandchars=\\\{\}}
% Add ',fontsize=\small' for more characters per line
\usepackage{framed}
\definecolor{shadecolor}{RGB}{241,243,245}
\newenvironment{Shaded}{\begin{snugshade}}{\end{snugshade}}
\newcommand{\AlertTok}[1]{\textcolor[rgb]{0.68,0.00,0.00}{#1}}
\newcommand{\AnnotationTok}[1]{\textcolor[rgb]{0.37,0.37,0.37}{#1}}
\newcommand{\AttributeTok}[1]{\textcolor[rgb]{0.40,0.45,0.13}{#1}}
\newcommand{\BaseNTok}[1]{\textcolor[rgb]{0.68,0.00,0.00}{#1}}
\newcommand{\BuiltInTok}[1]{\textcolor[rgb]{0.00,0.23,0.31}{#1}}
\newcommand{\CharTok}[1]{\textcolor[rgb]{0.13,0.47,0.30}{#1}}
\newcommand{\CommentTok}[1]{\textcolor[rgb]{0.37,0.37,0.37}{#1}}
\newcommand{\CommentVarTok}[1]{\textcolor[rgb]{0.37,0.37,0.37}{\textit{#1}}}
\newcommand{\ConstantTok}[1]{\textcolor[rgb]{0.56,0.35,0.01}{#1}}
\newcommand{\ControlFlowTok}[1]{\textcolor[rgb]{0.00,0.23,0.31}{\textbf{#1}}}
\newcommand{\DataTypeTok}[1]{\textcolor[rgb]{0.68,0.00,0.00}{#1}}
\newcommand{\DecValTok}[1]{\textcolor[rgb]{0.68,0.00,0.00}{#1}}
\newcommand{\DocumentationTok}[1]{\textcolor[rgb]{0.37,0.37,0.37}{\textit{#1}}}
\newcommand{\ErrorTok}[1]{\textcolor[rgb]{0.68,0.00,0.00}{#1}}
\newcommand{\ExtensionTok}[1]{\textcolor[rgb]{0.00,0.23,0.31}{#1}}
\newcommand{\FloatTok}[1]{\textcolor[rgb]{0.68,0.00,0.00}{#1}}
\newcommand{\FunctionTok}[1]{\textcolor[rgb]{0.28,0.35,0.67}{#1}}
\newcommand{\ImportTok}[1]{\textcolor[rgb]{0.00,0.46,0.62}{#1}}
\newcommand{\InformationTok}[1]{\textcolor[rgb]{0.37,0.37,0.37}{#1}}
\newcommand{\KeywordTok}[1]{\textcolor[rgb]{0.00,0.23,0.31}{\textbf{#1}}}
\newcommand{\NormalTok}[1]{\textcolor[rgb]{0.00,0.23,0.31}{#1}}
\newcommand{\OperatorTok}[1]{\textcolor[rgb]{0.37,0.37,0.37}{#1}}
\newcommand{\OtherTok}[1]{\textcolor[rgb]{0.00,0.23,0.31}{#1}}
\newcommand{\PreprocessorTok}[1]{\textcolor[rgb]{0.68,0.00,0.00}{#1}}
\newcommand{\RegionMarkerTok}[1]{\textcolor[rgb]{0.00,0.23,0.31}{#1}}
\newcommand{\SpecialCharTok}[1]{\textcolor[rgb]{0.37,0.37,0.37}{#1}}
\newcommand{\SpecialStringTok}[1]{\textcolor[rgb]{0.13,0.47,0.30}{#1}}
\newcommand{\StringTok}[1]{\textcolor[rgb]{0.13,0.47,0.30}{#1}}
\newcommand{\VariableTok}[1]{\textcolor[rgb]{0.07,0.07,0.07}{#1}}
\newcommand{\VerbatimStringTok}[1]{\textcolor[rgb]{0.13,0.47,0.30}{#1}}
\newcommand{\WarningTok}[1]{\textcolor[rgb]{0.37,0.37,0.37}{\textit{#1}}}

\providecommand{\tightlist}{%
  \setlength{\itemsep}{0pt}\setlength{\parskip}{0pt}}\usepackage{longtable,booktabs,array}
\usepackage{calc} % for calculating minipage widths
% Correct order of tables after \paragraph or \subparagraph
\usepackage{etoolbox}
\makeatletter
\patchcmd\longtable{\par}{\if@noskipsec\mbox{}\fi\par}{}{}
\makeatother
% Allow footnotes in longtable head/foot
\IfFileExists{footnotehyper.sty}{\usepackage{footnotehyper}}{\usepackage{footnote}}
\makesavenoteenv{longtable}
\usepackage{graphicx}
\makeatletter
\def\maxwidth{\ifdim\Gin@nat@width>\linewidth\linewidth\else\Gin@nat@width\fi}
\def\maxheight{\ifdim\Gin@nat@height>\textheight\textheight\else\Gin@nat@height\fi}
\makeatother
% Scale images if necessary, so that they will not overflow the page
% margins by default, and it is still possible to overwrite the defaults
% using explicit options in \includegraphics[width, height, ...]{}
\setkeys{Gin}{width=\maxwidth,height=\maxheight,keepaspectratio}
% Set default figure placement to htbp
\makeatletter
\def\fps@figure{htbp}
\makeatother

\KOMAoption{captions}{tableheading}
\makeatletter
\@ifpackageloaded{tcolorbox}{}{\usepackage[skins,breakable]{tcolorbox}}
\@ifpackageloaded{fontawesome5}{}{\usepackage{fontawesome5}}
\definecolor{quarto-callout-color}{HTML}{909090}
\definecolor{quarto-callout-note-color}{HTML}{0758E5}
\definecolor{quarto-callout-important-color}{HTML}{CC1914}
\definecolor{quarto-callout-warning-color}{HTML}{EB9113}
\definecolor{quarto-callout-tip-color}{HTML}{00A047}
\definecolor{quarto-callout-caution-color}{HTML}{FC5300}
\definecolor{quarto-callout-color-frame}{HTML}{acacac}
\definecolor{quarto-callout-note-color-frame}{HTML}{4582ec}
\definecolor{quarto-callout-important-color-frame}{HTML}{d9534f}
\definecolor{quarto-callout-warning-color-frame}{HTML}{f0ad4e}
\definecolor{quarto-callout-tip-color-frame}{HTML}{02b875}
\definecolor{quarto-callout-caution-color-frame}{HTML}{fd7e14}
\makeatother
\makeatletter
\@ifpackageloaded{caption}{}{\usepackage{caption}}
\AtBeginDocument{%
\ifdefined\contentsname
  \renewcommand*\contentsname{Table of contents}
\else
  \newcommand\contentsname{Table of contents}
\fi
\ifdefined\listfigurename
  \renewcommand*\listfigurename{List of Figures}
\else
  \newcommand\listfigurename{List of Figures}
\fi
\ifdefined\listtablename
  \renewcommand*\listtablename{List of Tables}
\else
  \newcommand\listtablename{List of Tables}
\fi
\ifdefined\figurename
  \renewcommand*\figurename{Figure}
\else
  \newcommand\figurename{Figure}
\fi
\ifdefined\tablename
  \renewcommand*\tablename{Table}
\else
  \newcommand\tablename{Table}
\fi
}
\@ifpackageloaded{float}{}{\usepackage{float}}
\floatstyle{ruled}
\@ifundefined{c@chapter}{\newfloat{codelisting}{h}{lop}}{\newfloat{codelisting}{h}{lop}[chapter]}
\floatname{codelisting}{Listing}
\newcommand*\listoflistings{\listof{codelisting}{List of Listings}}
\makeatother
\makeatletter
\makeatother
\makeatletter
\@ifpackageloaded{caption}{}{\usepackage{caption}}
\@ifpackageloaded{subcaption}{}{\usepackage{subcaption}}
\makeatother
\makeatletter
\@ifpackageloaded{fontawesome5}{}{\usepackage{fontawesome5}}
\makeatother

\ifLuaTeX
  \usepackage{selnolig}  % disable illegal ligatures
\fi
\usepackage{bookmark}

\IfFileExists{xurl.sty}{\usepackage{xurl}}{} % add URL line breaks if available
\urlstyle{same} % disable monospaced font for URLs
\hypersetup{
  pdftitle={Ayman's Workshop Series},
  pdfauthor={Ayman Yousef},
  colorlinks=true,
  linkcolor={blue},
  filecolor={Maroon},
  citecolor={Blue},
  urlcolor={Blue},
  pdfcreator={LaTeX via pandoc}}


\title{Ayman's Workshop Series}
\author{Ayman Yousef}
\date{2024-09-13}

\begin{document}
\maketitle


\section{}\label{section}

{Workshop 1: Introduction to the Randles Lab}\\
{Getting started with our different tools}\\
{~}\\
\href{https://aymanzyy.github.io}{\faIcon{home}} {Ayman Yousef }

{2024-09-04 @ \href{https://randleslab.pratt.duke.edu}{Randles Lab
Workshop Series}}

\section{Overview}\label{overview}

\begin{enumerate}
\def\labelenumi{\arabic{enumi}.}
\tightlist
\item
  \hyperref[background]{Background}

  \begin{itemize}
  \tightlist
  \item
    \hyperref[wecare]{Why Do We Care}
  \end{itemize}
\item
  \hyperref[documentation]{Documentation}

  \begin{itemize}
  \tightlist
  \item
    \hyperref[wiki]{Wiki}
  \item
    \hyperref[box]{Box}
  \item
    \hyperref[outlook]{Outlook}
  \end{itemize}
\item
  \hyperref[monday]{Monday.com}
\item
  \hyperref[overleaf]{Overleaf}
\item
  \hyperref[whitesides]{Whiteside's}
\end{enumerate}

\section{Background}\label{background}

\begin{tcolorbox}[enhanced jigsaw, colframe=quarto-callout-note-color-frame, bottomrule=.15mm, toptitle=1mm, opacitybacktitle=0.6, colback=white, opacityback=0, colbacktitle=quarto-callout-note-color!10!white, titlerule=0mm, toprule=.15mm, coltitle=black, left=2mm, arc=.35mm, bottomtitle=1mm, rightrule=.15mm, title=\textcolor{quarto-callout-note-color}{\faInfo}\hspace{0.5em}{Goal}, breakable, leftrule=.75mm]

Get you guys caught up on the different resources available in the lab.

\end{tcolorbox}

\begin{itemize}
\item
  It's \textbf{important} to note that the expectation is to keep up
  with each of these resources.
\item
  {Ex}: Monday pages are checked at random and during project meetings
\end{itemize}

\includegraphics[width=0.3\textwidth,height=\textheight]{index_files/mediabag/confluence-logo-D9B0.png}
\includegraphics[width=0.5\textwidth,height=\textheight]{index_files/mediabag/Logo-1-NO-1.png}

You are not expected to be experts by the end of this, but familiarty
with each of these is {expected}.

\section{Why Do We Care}\label{wecare}

\begin{itemize}
\tightlist
\item
  Reproducibility is one of the most important aspets of research
\item
  Key tenant of the Randles Lab, you will be involved in multiple
  projects
\end{itemize}

\begin{tcolorbox}[enhanced jigsaw, colframe=quarto-callout-warning-color-frame, bottomrule=.15mm, toptitle=1mm, opacitybacktitle=0.6, colback=white, opacityback=0, colbacktitle=quarto-callout-warning-color!10!white, titlerule=0mm, toprule=.15mm, coltitle=black, left=2mm, arc=.35mm, bottomtitle=1mm, rightrule=.15mm, title=\textcolor{quarto-callout-warning-color}{\faExclamationTriangle}\hspace{0.5em}{Warning!}, breakable, leftrule=.75mm]

\begin{itemize}
\tightlist
\item
  Your focus will be diverted on a day-to-day basis. Good to make notes
  for both yourself and fellow lab members
\end{itemize}

\end{tcolorbox}

\begin{itemize}
\tightlist
\item
  Collaboration is a key aspect of our research. Courteous to keep good
  notes so others can easily replicate your work

  \begin{itemize}
  \tightlist
  \item
    {Ex}: Getting hit by a bus
  \end{itemize}
\end{itemize}

\begin{itemize}
\tightlist
\item
  Learning these tools are important in sustaining a collaborative lab
  environment.

  \begin{itemize}
  \tightlist
  \item
    If we're all comfortable with the same tools, it becomes easier pass
    of tasks, share work, etc.
  \end{itemize}
\end{itemize}

\section{Documentation}\label{documentation}

\begin{itemize}
\item
  Wiki (Confluence)
\item
  Box
\item
  Outlook
\end{itemize}

\section{Wiki (Confluence)}\label{wiki}

\begin{itemize}
\item
  Website that houses most of our documentation
\item
  Project updates are written here, with everyone in the lab having
  access to every page
\item
  \textbf{{Structure}}:

  \begin{itemize}
  \tightlist
  \item
    Each project has a ``landing page'' and a ``notebook''
  \item
    Notebooks contain information regarding the ongoing project it's for
  \end{itemize}
\end{itemize}

\includegraphics[width=0.9\textwidth,height=\textheight]{index_files/mediabag/66b9f922cf03aa041292.png}

\includegraphics[width=0.9\textwidth,height=\textheight]{index_files/mediabag/48c5538a154172f98d50.png}

\subsection{Wiki: Sample Entry}\label{wiki-sample-entry}

\begin{itemize}
\tightlist
\item
  {\textbf{Date} }:

  \begin{itemize}
  \tightlist
  \item
    Can be days when work was performed or when journal entry is
    recorded
  \end{itemize}
\item
  {\textbf{Summary} }:

  \begin{itemize}
  \tightlist
  \item
    Include context for what work you're performing.
  \item
    Can be conversational, some of the best notebooks work through the
    thinking process that led to a specific run, solution, etc.
  \end{itemize}
\end{itemize}

\begin{tcolorbox}[enhanced jigsaw, colframe=quarto-callout-color-frame, bottomrule=.15mm, toprule=.15mm, left=2mm, arc=.35mm, colback=white, rightrule=.15mm, opacityback=0, breakable, leftrule=.75mm]

\vspace{-3mm}\textbf{Algorithm}\vspace{3mm}

\begin{enumerate}
\def\labelenumi{\arabic{enumi}.}
\item
  \texttt{Date}:
\item
  \texttt{Goal}:
\item
  \texttt{Input\ File}:
\item
  \texttt{Path\ to\ data\ output}:
\item
  \texttt{Code\ Version\ -\ git\ checkout}:

  \begin{itemize}
  \tightlist
  \item
    Link to branch
  \item
    Write commit hash
  \end{itemize}
\item
  \texttt{Job\ Submission\ Parameters}:

  \begin{itemize}
  \tightlist
  \item
    Number of processors
  \item
    Relevant job parameters

    \begin{itemize}
    \tightlist
    \item
      Partition
    \item
      Node numbers
    \item
      Flags
    \end{itemize}
  \end{itemize}
\item
  \texttt{Outcome/Post\ Processing}:

  \begin{itemize}
  \tightlist
  \item
    Link scripts
  \item
    Explain your workflow
  \end{itemize}
\end{enumerate}

\end{tcolorbox}

\subsection{Wiki: Example Entry}\label{wiki-example-entry}

\includegraphics{index_files/mediabag/4315a92dce655ed54b03.png}

\subsection{Wiki: Individual Meeting
Sign-Up}\label{wiki-individual-meeting-sign-up}

\includegraphics[width=0.6\textwidth,height=\textheight]{index_files/mediabag/aa354f6cbe0674f3380a.png}

\subsection{Wiki: Group Meeting
Schedule}\label{wiki-group-meeting-schedule}

\includegraphics{index_files/mediabag/95a2c75fde6d2d5066f4.png}

\subsection{Wiki: Lattice Boltzmann
Questions}\label{wiki-lattice-boltzmann-questions}

\begin{itemize}
\item
  Since not all of us have an extensive fluid background, these
  questions help catch everyone up
\item
  Questions are split into two parts, one part done you first semester
  and the other done your second semester
\end{itemize}

\includegraphics{index_files/mediabag/d5e7fa20ba84dd451270.png}

\section{Box}\label{box}

\begin{itemize}
\tightlist
\item
  Website that houses the rest of our documentation
\item
  Files to be shared with the lab, Professor Randles are housed here

  \begin{itemize}
  \tightlist
  \item
    Group Presentations
  \item
    Individual meeting files
  \item
    Milestone documents
  \end{itemize}
\item
  Unlimited lab storage but Duke also gives individuals their own
  allocation of storage
\end{itemize}

\begin{tcolorbox}[enhanced jigsaw, colframe=quarto-callout-warning-color-frame, bottomrule=.15mm, toptitle=1mm, opacitybacktitle=0.6, colback=white, opacityback=0, colbacktitle=quarto-callout-warning-color!10!white, titlerule=0mm, toprule=.15mm, coltitle=black, left=2mm, arc=.35mm, bottomtitle=1mm, rightrule=.15mm, title=\textcolor{quarto-callout-warning-color}{\faExclamationTriangle}\hspace{0.5em}{Warning!}, breakable, leftrule=.75mm]

\begin{itemize}
\tightlist
\item
  DO NOT UPLOAD UNNECESSARY FILES
\end{itemize}

\end{tcolorbox}

\section{Outlook}\label{outlook}

\begin{itemize}
\tightlist
\item
  We use Outlook calendar to organize our lab schedule. This includes:

  \begin{itemize}
  \tightlist
  \item
    Meetings
  \item
    Lab Activities
  \item
    Milestone Events
  \item
    Time Off
  \end{itemize}
\end{itemize}

\begin{tcolorbox}[enhanced jigsaw, colframe=quarto-callout-important-color-frame, bottomrule=.15mm, toptitle=1mm, opacitybacktitle=0.6, colback=white, opacityback=0, colbacktitle=quarto-callout-important-color!10!white, titlerule=0mm, toprule=.15mm, coltitle=black, left=2mm, arc=.35mm, bottomtitle=1mm, rightrule=.15mm, title=\textcolor{quarto-callout-important-color}{\faExclamation}\hspace{0.5em}{Warning}, breakable, leftrule=.75mm]

\textbf{Professor Randles lives by her calendar. If it's not on her
calendar, it's not happening}:

\end{tcolorbox}

\includegraphics[width=0.75\textwidth,height=\textheight]{index_files/mediabag/Microsoft_Office_Out.png}

\subsection{Outlook Overview}\label{outlook-overview}

\includegraphics[width=0.75\textwidth,height=\textheight]{index_files/mediabag/cadbca20725120ffc1d6.png}

\subsection{Outlook Day-to-Day}\label{outlook-day-to-day}

\includegraphics[width=0.65\textwidth,height=\textheight]{index_files/mediabag/cbc94a3f77ad59c33898.png}

\section{Monday.com}\label{monday}

\begin{itemize}
\tightlist
\item
  Monday.com Ltd.~is a cloud-based platform that allows users to create
  their own applications and project management software.
\item
  It is used to schedule/assign tasks, track project progress, etc.
\end{itemize}

\begin{tcolorbox}[enhanced jigsaw, colframe=quarto-callout-warning-color-frame, bottomrule=.15mm, toptitle=1mm, opacitybacktitle=0.6, colback=white, opacityback=0, colbacktitle=quarto-callout-warning-color!10!white, titlerule=0mm, toprule=.15mm, coltitle=black, left=2mm, arc=.35mm, bottomtitle=1mm, rightrule=.15mm, title=\textcolor{quarto-callout-warning-color}{\faExclamationTriangle}\hspace{0.5em}{Warning!}, breakable, leftrule=.75mm]

\begin{itemize}
\tightlist
\item
  \textbf{DO NOT UPLOAD UNNECESSARY FILES}
\item
  \textbf{DO NOT UPLOAD PATIENT DATA}
\end{itemize}

\end{tcolorbox}

\includegraphics[width=0.75\textwidth,height=\textheight]{index_files/mediabag/In-depth-review-–-.png}

\subsection{Monday.com: Organizational
Structure}\label{monday.com-organizational-structure}

\begin{itemize}
\tightlist
\item
  {\textbf{My Work} }:

  \begin{itemize}
  \tightlist
  \item
    View all tasks assigned to yourself across all boards
  \item
    Customize people, dates, and status columns and choose which boards
    to view
  \end{itemize}
\item
  {\textbf{Meetings} }

  \begin{itemize}
  \tightlist
  \item
    Group Meetings
  \item
    Project meetings
  \end{itemize}
\item
  { \textbf{Projects} }

  \begin{itemize}
  \tightlist
  \item
    In each project group folder, there is a dashboard that collects
    info from all projects and a separate board for each project
  \end{itemize}
\end{itemize}

\begin{itemize}
\tightlist
\item
  {\textbf{Resources} }

  \begin{itemize}
  \tightlist
  \item
    Instructions/Tutorials for using Monday
  \item
    Contacts
  \item
    Project List
  \end{itemize}
\item
  {\textbf{Miscellaneous Work} }

  \begin{itemize}
  \tightlist
  \item
    Monday (requests and Monday maintenance)
  \item
    Miscellaneous Task Board
  \item
    Journals
  \end{itemize}
\end{itemize}

\subsection{Monday.com: My Work Page}\label{monday.com-my-work-page}

\includegraphics[width=0.75\textwidth,height=\textheight]{index_files/mediabag/5c9160cabb7ca36325d9.png}

\subsection{Monday.com: Project Page}\label{monday.com-project-page}

\includegraphics[width=0.75\textwidth,height=\textheight]{index_files/mediabag/9d73bb314254e474d7fa.png}

\subsection{Monday.com: Research Subgroup
Page}\label{monday.com-research-subgroup-page}

\includegraphics[width=0.75\textwidth,height=\textheight]{index_files/mediabag/70884729c5230109d0e2.png}

\subsection{Monday.com: Board Organization
{[}Views{]}}\label{monday.com-board-organization-views}

\begin{itemize}
\tightlist
\item
  {\textbf{Project} }

  \begin{itemize}
  \tightlist
  \item
    This should be the default view for all projects\\
  \item
    Includes all critical columns
  \end{itemize}
\item
  {\textbf{Main Table} }

  \begin{itemize}
  \tightlist
  \item
    In general is not used
  \item
    Includes every column that can be used, including back-end triggers
    (used to trigger automations) which users should not edit
  \end{itemize}
\end{itemize}

\begin{itemize}
\tightlist
\item
  {\textbf{Papers} }

  \begin{itemize}
  \tightlist
  \item
    includes columns that are only relevant to papers (Other authors,
    paper stage, link to whitesides, Conference submitting to, etc.)
  \end{itemize}
\item
  {\textbf{Instructions} }

  \begin{itemize}
  \tightlist
  \item
    View that you should see when the board is initialized
  \item
    Contains instructions for properly setting up the project
  \item
    Tasks must be completed in a timely fashion
  \end{itemize}
\end{itemize}

\subsection{Monday.com: Board Organization
{[}Columns{]}}\label{monday.com-board-organization-columns}

\begin{itemize}
\tightlist
\item
  {\textbf{Person} }

  \begin{itemize}
  \tightlist
  \item
    Individual(s) who are assigned to the task
  \item
    Used by My Work to identify tasks
  \end{itemize}
\item
  {\textbf{Status} }

  \begin{itemize}
  \tightlist
  \item
    Indicates if the task is done, stuck, working on it, etc.
  \item
    Done vs Done and Discussed
  \end{itemize}
\item
  {\textbf{Dependent On} }

  \begin{itemize}
  \tightlist
  \item
    Used to set a difference in deadlines between two or more tasks.
  \item
    Used for all paper templates to set deadlines based on submission
    date
  \end{itemize}
\end{itemize}

\begin{itemize}
\tightlist
\item
  {\textbf{Due Date} }

  \begin{itemize}
  \tightlist
  \item
    Used to organize tasks by date in My Work
  \end{itemize}
\item
  {\textbf{Link} }

  \begin{itemize}
  \tightlist
  \item
    Any useful links, used in the ``Wiki Link'' group to connect the
    project's notebook
  \end{itemize}
\item
  {\textbf{Project Leader} }

  \begin{itemize}
  \tightlist
  \item
    Should be set as a default value for the board
  \item
    Can be filtered for in My Work to see all tasks that they are
    supervising
  \end{itemize}
\item
  {\textbf{Notifications} }

  \begin{itemize}
  \tightlist
  \item
    Receive a notification 1 week, 3 days, or 1 day before a task is due
  \end{itemize}
\end{itemize}

\section{Overleaf}\label{overleaf}

\begin{itemize}
\tightlist
\item
  Overleaf ``is a collaborative cloud-based LaTeX editor used for
  writing, editing and publishing scientific documents''.
\item
  All papers, outlines, milestone documents, and some of our
  documentation is written with Overleaf. Easily shareable and great for
  bookkeeping
\end{itemize}

\begin{tcolorbox}[enhanced jigsaw, colframe=quarto-callout-tip-color-frame, bottomrule=.15mm, toptitle=1mm, opacitybacktitle=0.6, colback=white, opacityback=0, colbacktitle=quarto-callout-tip-color!10!white, titlerule=0mm, toprule=.15mm, coltitle=black, left=2mm, arc=.35mm, bottomtitle=1mm, rightrule=.15mm, title=\textcolor{quarto-callout-tip-color}{\faLightbulb}\hspace{0.5em}{Learning Overleaf!}, breakable, leftrule=.75mm]

\begin{itemize}
\tightlist
\item
  Learning LaTeX while your research expectations are lower is a good
  idea, trust me
\end{itemize}

\end{tcolorbox}

\subsection{Overleaf: LaTeX}\label{overleaf-latex}

\begin{itemize}
\tightlist
\item
  LATEX is a tool for typesetting professional-looking documents
\item
  LaTeX works very differently: instead of a interactive page to type
  and edit text with, your document is a plain text file interspersed
  with LaTeX commands used to express the desired (typeset) results.
\item
  To produce a visible, typeset document, your LaTeX file is processed
  by a piece of software called a TeX engine which uses the commands
  embedded in your text file to guide and control the typesetting
  process, converting the LaTeX commands and document text into a
  professionally typeset PDF file
\end{itemize}

\subsection{Overleaf: Writing your first piece of
LATEX}\label{overleaf-writing-your-first-piece-of-latex}

\begin{itemize}
\item ~
  \subsection{Creating a document and writing some
  text}\label{creating-a-document-and-writing-some-text}

  \begin{codelisting}

  \caption{\texttt{Document Creation}}

\begin{Shaded}
\begin{Highlighting}[numbers=left,,]
\BuiltInTok{\textbackslash{}documentclass}\NormalTok{\{}\ExtensionTok{article}\NormalTok{\}}
\KeywordTok{\textbackslash{}begin}\NormalTok{\{}\ExtensionTok{document}\NormalTok{\}}
\NormalTok{First document. This is a simple example, with no }
\NormalTok{extra parameters or packages included.}
\KeywordTok{\textbackslash{}end}\NormalTok{\{}\ExtensionTok{document}\NormalTok{\}}
\end{Highlighting}
\end{Shaded}

  \end{codelisting}

  \begin{itemize}
  \tightlist
  \item
    Will produce the following image

    \begin{itemize}
    \tightlist
    \item
      \includegraphics[width=0.75\textwidth,height=\textheight]{index_files/mediabag/9b604fe70b5c218c6bc6.png}
    \end{itemize}
  \end{itemize}
\item
  Not always this easy. Some journals/conferences require specific
  formats
\end{itemize}

\subsection{Overleaf: Document
Pre-amble}\label{overleaf-document-pre-amble}

\begin{itemize}
\tightlist
\item
  Everything in your .tex file appearing before \textbackslash begin
  document is called the preamble, which acts as the document's
  ``setup'' section.
\end{itemize}

\begin{codelisting}

\caption{\texttt{Document Settings}}

\begin{Shaded}
\begin{Highlighting}[]
\BuiltInTok{\textbackslash{}documentclass}\NormalTok{[12pt, letterpaper]\{}\ExtensionTok{article}\NormalTok{\}}
\BuiltInTok{\textbackslash{}usepackage}\NormalTok{\{}\ExtensionTok{graphicx}\NormalTok{\}}
\end{Highlighting}
\end{Shaded}

\end{codelisting}

\begin{itemize}
\tightlist
\item
  In this example, we are setting both:

  \begin{itemize}
  \tightlist
  \item
    The {font size} to 12 pt
  \item
    The {paper size} to letterpaper
  \end{itemize}
\end{itemize}

\begin{itemize}
\tightlist
\item
  Overleaf is compatible with a bunch of different libraries, packages
  that you can use to your advantage. To do so, you'd want to include
  something like the following:
\end{itemize}

\begin{codelisting}

\caption{\texttt{Package Import}}

\begin{Shaded}
\begin{Highlighting}[]
\BuiltInTok{\textbackslash{}usepackage}\NormalTok{\{}\ExtensionTok{graphicx}\NormalTok{\}}
\end{Highlighting}
\end{Shaded}

\end{codelisting}

\begin{itemize}
\tightlist
\item
  For more info on what packages are available and their functionality,
  you can visit the following
  \href{https://www.overleaf.com/learn/latex/Learn_LaTeX_in_30_minutes\#Finding_and_using_LaTeX_packages}{page}
\end{itemize}

\subsection{Overleaf: Including title, author and date
information}\label{overleaf-including-title-author-and-date-information}

\begin{itemize}
\tightlist
\item
  Adding a title, author and date to our document requires three more
  lines in the preamble (not the main body of the document). Those lines
  are:
\end{itemize}

\begin{Shaded}
\begin{Highlighting}[]
\FunctionTok{\textbackslash{}title}\NormalTok{\{My first LaTeX document\}: }
\FunctionTok{\textbackslash{}author}\NormalTok{\{Hubert Farnsworth\}: }

\FunctionTok{\textbackslash{}thanks}\NormalTok{\{Funded by the Overleaf team.\}: }

\FunctionTok{\textbackslash{}date}\NormalTok{\{August 2022\}: }
\end{Highlighting}
\end{Shaded}

\subsection{Overleaf: Adding comments}\label{overleaf-adding-comments}

To make a comment in LATEX, simply write a {\%} symbol at the beginning
of the line, as shown in the following code which uses the example
above:

\begin{codelisting}

\caption{\texttt{Example Comment Line}}

\begin{Shaded}
\begin{Highlighting}[]
\CommentTok{\% This line here is a comment. It will not be typeset in the document.}
\end{Highlighting}
\end{Shaded}

\end{codelisting}

\subsection{Overleaf: Bold, italics and
underlining}\label{overleaf-bold-italics-and-underlining}

\textbf{Bold}: bold text in LaTeX is typeset using the \textbf{...}
command.

\begin{codelisting}

\caption{\texttt{Example of Bolding}}

\begin{Shaded}
\begin{Highlighting}[]
\NormalTok{Some of the }\FunctionTok{\textbackslash{}textbf}\NormalTok{\{greatest\}}
\end{Highlighting}
\end{Shaded}

\end{codelisting}

Underlined: to underline text use the \underline{...} command.

\begin{codelisting}

\caption{\texttt{Example of Underlining}}

\begin{Shaded}
\begin{Highlighting}[]
\NormalTok{discoveries in }\FunctionTok{\textbackslash{}underline}\NormalTok{\{science\} }
\end{Highlighting}
\end{Shaded}

\end{codelisting}

\emph{Italics}: italicised text is produced using the \textit{...}
command.

\begin{codelisting}

\caption{\texttt{Example of Italicizing}}

\begin{Shaded}
\begin{Highlighting}[]
\NormalTok{were made by }\FunctionTok{\textbackslash{}textbf}\NormalTok{\{}\FunctionTok{\textbackslash{}textit}\NormalTok{\{accident\}\}.}
\end{Highlighting}
\end{Shaded}

\end{codelisting}

\subsection{Overleaf: Adding Images}\label{overleaf-adding-images}

Overleaf supports three ways to insert images:

\begin{itemize}
\tightlist
\item
  Use the \textbf{Insert Figure} button, located on the editor toolbar,
  to insert an image into Visual Editor or Code Editor.
\item
  Copy and paste an image into Visual Editor or Code Editor.
\item
  Use Code Editor to write LaTeX code that inserts a graphic.
\end{itemize}

\begin{Shaded}
\begin{Highlighting}[numbers=left,,]
\BuiltInTok{\textbackslash{}documentclass}\NormalTok{\{}\ExtensionTok{article}\NormalTok{\}}
\BuiltInTok{\textbackslash{}usepackage}\NormalTok{\{}\ExtensionTok{graphicx}\NormalTok{\} }\CommentTok{\%LaTeX package to import graphics}
\FunctionTok{\textbackslash{}graphicspath}\NormalTok{\{\{images/\}\} }\CommentTok{\%configuring the graphicx package}
 
\KeywordTok{\textbackslash{}begin}\NormalTok{\{}\ExtensionTok{document}\NormalTok{\}}
\NormalTok{The universe is immense and it seems to be homogeneous, }
\NormalTok{on a large scale, everywhere we look.}

\CommentTok{\% The \textbackslash{}includegraphcs command is }
\CommentTok{\% provided (implemented) by the }
\CommentTok{\% graphicx package}
\BuiltInTok{\textbackslash{}includegraphics}\NormalTok{\{}\ExtensionTok{universe}\NormalTok{\}  }
 
\NormalTok{There\textquotesingle{}s a picture of a galaxy above (trust me yall).}
\KeywordTok{\textbackslash{}end}\NormalTok{\{}\ExtensionTok{document}\NormalTok{\}}
\end{Highlighting}
\end{Shaded}

\subsection{Overleaf: Captions, labels,
references}\label{overleaf-captions-labels-references}

\begin{Shaded}
\begin{Highlighting}[numbers=left,,]
\BuiltInTok{\textbackslash{}documentclass}\NormalTok{\{}\ExtensionTok{article}\NormalTok{\}}
\BuiltInTok{\textbackslash{}usepackage}\NormalTok{\{}\ExtensionTok{graphicx}\NormalTok{\}}
\FunctionTok{\textbackslash{}graphicspath}\NormalTok{\{\{images/\}\}}
 
\KeywordTok{\textbackslash{}begin}\NormalTok{\{}\ExtensionTok{document}\NormalTok{\}}

\KeywordTok{\textbackslash{}begin}\NormalTok{\{}\ExtensionTok{figure}\NormalTok{\}[h]}
    \FunctionTok{\textbackslash{}centering}
    \BuiltInTok{\textbackslash{}includegraphics}\NormalTok{[width=0.75}\FunctionTok{\textbackslash{}textwidth}\NormalTok{]\{}\ExtensionTok{mesh}\NormalTok{\}}
    \FunctionTok{\textbackslash{}caption}\NormalTok{\{A nice plot.\}}
    \KeywordTok{\textbackslash{}label}\NormalTok{\{}\ExtensionTok{fig:mesh1}\NormalTok{\}}
\KeywordTok{\textbackslash{}end}\NormalTok{\{}\ExtensionTok{figure}\NormalTok{\}}
 
\NormalTok{As you can see in figure }\KeywordTok{\textbackslash{}ref}\NormalTok{\{}\ExtensionTok{fig:mesh1}\NormalTok{\}, the function grows near the origin. This example is on page }\KeywordTok{\textbackslash{}pageref}\NormalTok{\{}\ExtensionTok{fig:mesh1}\NormalTok{\}.}

\KeywordTok{\textbackslash{}end}\NormalTok{\{}\ExtensionTok{document}\NormalTok{\}}
\end{Highlighting}
\end{Shaded}

\subsection{Overleaf: Link Sharing}\label{overleaf-link-sharing}

\subsection{Overleaf: More to Learn}\label{overleaf-more-to-learn}

\textbf{There's so much more I can cover, but I won't}

\begin{itemize}
\tightlist
\item
  Creating lists
\item
  Adding math
\item
  Basic document structure
\item
  Creating tables
\item
  Community formats
\item
  LaTeX packages
\end{itemize}

Useful Resources include:

\begin{itemize}
\item
  Overleaf's Website
\item
  Dr Trefor Bazett's Overleaf Series

  \begin{itemize}
  \item
    \url{https://www.youtube.com/watch?v=-HvRvBjBAvg}
  \end{itemize}
\item
  Everyone in the lab (they should know these things)
\end{itemize}

\section{Whiteside's Outline}\label{whitesides-outline}

\begin{itemize}
\tightlist
\item
  Based on a paper by Harvard Professor/American chemist George
  Whiteside
\item
  Idea behind his work on papers was that ``it is most efficient to
  write papers from outlines''
\item
  You can find the template
  \href{https://duke.app.box.com/file/92773425392}{here}
\end{itemize}

\includegraphics[width=0.75\textwidth,height=\textheight]{index_files/mediabag/800wm.pdf}

\subsection{Whiteside's Outline:
Overview}\label{whitesides-outline-overview}

I'm not teaching you this, you gotta look through the documentation in
the Randles Lab Document




\end{document}
